\chapter{Discussion}
\label{ch:discussion}

The results of Chapter~\ref{ch:results} invert the expectation that frames most of
the cooperative-MARL literature. MAPPO is the sophisticated, theoretically motivated
choice: its centralized critic is designed precisely to tame non-stationarity and to
solve the multi-agent credit-assignment problem, and its neighbour-coupled reward
explicitly encodes coordination. One naturally expects it to beat a set of
independent learners that ignore all of this. Yet on the \grid{2} network the two
are statistically indistinguishable on every performance metric, and the only
reliable difference --- MAPPO's slightly higher switching rate --- is, if anything, a
mild inefficiency rather than a benefit. This chapter asks why.

\section{Why the Coordination Value Is Near Zero}
The most parsimonious explanation is that, at this scale, \emph{there is little
coordination left for an explicit mechanism to add}. A \grid{2} grid has only four
intersections, each adjacent to two neighbours, and the traffic-control reward is
already locally informative: minimizing one's own queue and waiting time, on a
network this tightly coupled, is almost the same objective as minimizing the
network's. In other words, the incentives are \emph{already aligned} by the
structure of the problem. The centralized critic provides a lower-variance value
estimate and the neighbour-pressure term nudges toward network-aware behaviour, but
an independent learner equipped with the same local observations and the same
parameter-sharing can discover the same coordinated phasing on its own. Coordination
here is \emph{emergent} rather than \emph{engineered}: it falls out of independent
optimization because the environment does not reward defection.

This reading is reinforced by two details. First, both learned controllers
\emph{decisively} beat the heuristic baselines under matched clearance
(Chapter~\ref{ch:results}); the learning is doing real work --- it is not that the
task is trivial, but that the \emph{coordination-specific} additions are redundant.
Second, MAPPO's extra phase switches buy no improvement, which is consistent with a
centralized critic that is active but not pivotal: it shapes the policy slightly
differently without changing what the policy ultimately achieves.

\section{The Scale Hypothesis}
If the coordination value is small \emph{because} the network is small, then it
should grow with scale. As a grid enlarges, an agent's actions influence ever more
distant intersections through longer chains of traffic; the gap between ``minimize
my own delay'' and ``minimize the network's delay'' widens, non-stationarity from
the perspective of any single learner intensifies, and the centralized critic's
access to the global state should translate into a measurable advantage. This is
exactly the regime in which network-level cooperation has been shown to
matter: CoLight~\cite{wei2019colight} demonstrates gains from sharing neighbour
information at city scale, and the large-scale study of
Chu~et~al.~\cite{chu2020largescale} reports that neighbour-aware, carefully designed
shared rewards improve performance --- while naive individual rewards can actively
harm throughput --- on large networks. The near-zero coordination value measured here
is therefore not in tension with that literature; it is the expected
\emph{small-scale} end of the same curve.

\section{Implications}
The practical implication is a note of caution and a research direction. The caution
is that, on small or tightly coupled multi-agent problems, the considerable
machinery of CTDE may be unnecessary: a simpler independent-learning approach with
parameter sharing can match it, at lower implementation and training cost. The
research direction is that the \emph{value of coordination} should be treated as a
quantity that depends on system scale and on how strongly incentives conflict ---
not as a fixed property of an algorithm. The controlled traffic setting used here
establishes a clean zero-point for that quantity, against which both larger networks
and genuine social dilemmas (Chapter~\ref{ch:futurework}) can be measured.
